\section{The regular quotient and its intrinsic metric}
\label{sec:local-condition}
The global certificate $\eta_d$ and an intrinsic local condition answer
different questions. The former controls competitors throughout the
closed model. The latter is the norm of the differential of the target
on the observation image. We first identify that image, including all
its nearby closed-model representatives.

Let $\Theta^\circ_d$ consist of the parameters in $\mathcal R_d$ with
strictly positive channel entries, $\alpha_b>\alpha_*$, pairwise distinct
complete component polynomials, and simple component roots in $(L,D)$.
Use the first $m_i-1$ entries of each stochastic column, the first $k-1$
weights, and the nonleading coefficients of the $f_b$ as coordinates.
Thus $\Theta^\circ_d$ is an open subset of $\R^p$, where
\[
 p=k(m_1+m_2-2)+(k-1)+kd.
\]
The simple real-rooted coefficient region is open, by the implicit
function theorem and the sign-bracket argument. All rank and strict
inequality conditions above are open as well. The group $\mathfrak S_k$
acts by simultaneous permutation of components. Put
$F(\theta)=(\operatorname{vec}P_\theta(T_j))_j$.

\begin{theorem}[Local quotient identification]\label{thm:quotient-chart}
The action of $\mathfrak S_k$ on $\Theta^\circ_d$ is free. For every
$\theta_0\in\Theta^\circ_d$, with $P=F(\theta_0)$, there are an open
observation neighbourhood $\mathcal U$ of $P$ and a coordinate
neighbourhood $\mathcal V$ of $\theta_0$ such that:
\begin{enumerate}
\item every closed-model representative of a datum in
$\mathcal U\cap F(\cK_d)$ belongs to exactly one of the disjoint sets
$\pi\mathcal V$, $\pi\in\mathfrak S_k$;
\item each such datum has exactly one representative in $\mathcal V$;
\item the inverse from this model image to $\mathcal V$ is real analytic
and extends to a real-analytic map on an ambient neighbourhood of $P$.
\end{enumerate}
Consequently $\Theta^\circ_d/\mathfrak S_k$ is a manifold, not an
orbifold with nontrivial stabilizers, and its observation image is
locally an embedded analytic manifold. The derivative $J=DF(\theta_0)$
has full column rank. Collisions of roots belonging to different
component polynomials do not alter any of these conclusions.
\end{theorem}
\begin{proof}
\emph{Exact fibres.} A permutation fixing a parameter fixes its complete
polynomial tuple. Distinctness of these tuples makes the permutation
the identity. Suppose any parameter in the closed model represents $P$.
Since $\tau(P)>0$, its normalizing polynomial is the unique monic
solution of $\cN_Pa=b_P$. Interpolation therefore recovers the same
$K$. Its leading coefficient $A$ has rank $k$. The factorization
$A=U'\diag(\alpha')V'^{\mathsf T}$ forces both competing channels to
have full column rank. In fixed orthonormal frames the observable
coefficient algebra has $k$ distinct joint tuples, with rank-one joint
projectors. Any competing diagonal realization must assign exactly one
component to each tuple. For the projector $E_b$ of that tuple,
\begin{equation}\label{eq:rank-one-recovery}
 Q_b=Q_1E_bA_0Q_2^{\mathsf T}
       =\alpha_bU_{:b}V_{:b}^{\mathsf T}.
\end{equation}
Its entry sum is $\alpha_b$; its row and column sums, divided by
$\alpha_b$, are the two stochastic columns. Thus the entire exact fibre
in $\cK_d$ is the single permutation orbit of $\theta_0$.

\emph{Analytic recovery.} For an ambient datum $Y$ near $P$, define
\[
 a(Y)=(\cN_Y^{\mathsf T}\cN_Y)^{-1}\cN_Y^{\mathsf T}b_Y,
 \qquad K(Y)=I_d(q_Y Y).
\]
The least-squares inverse exists on a neighbourhood of $P$. Select $k$
independent columns and $k$ independent rows of the leading coefficient
$A(P)$. The same selections in $A(Y)$, orthonormalized by the analytic
positive definite inverse square root of their Gram matrices, give
analytic frames $Q_1(Y),Q_2(Y)$. On rank-$k$ model data they span the
whole column and row spaces. The reduced leading matrix $A_0(Y)$ is
invertible near $P$, so each $C_j(Y)=M_j(Y)A_0(Y)^{-1}$ is analytic.
Choose fixed real coefficients $c_j$ for which $\sum c_j C_j(P)$ has
$k$ distinct eigenvalues. Disjoint small circles about these eigenvalues
define analytic rank-one Riesz projectors $E_b(Y)$; simple real
eigenvalues remain real on this real neighbourhood. Recover the
coefficient tuple by $f_{bj}(Y)=\operatorname{tr}(C_j(Y)E_b(Y))$, and
recover the factors from \eqref{eq:rank-one-recovery} and its row,
column and total sums. These formulae are analytic ambient maps.
For data on the model, they give exactly its parameters, with the
labelling selected by the circles.

\emph{All nearby representatives.} The recovered parameters tend to
$\theta_0$ as $Y\to P$. Shrink the neighbourhood so that the recovered
weights, entries, ranks, distinct tuples and simple interior roots
satisfy the strict conditions defining $\Theta^\circ_d$. Every
closed-model representative of such a datum has the recovered $q,K$.
The rank argument and the singleton projectors just used apply to that
datum itself, not only to its limit. Its representatives are therefore
precisely the permutations of the recovered point. Choose $\mathcal V$
small enough that its finitely many permutation translates are disjoint,
and shrink $\mathcal U$ once more. This proves the first two assertions
without assuming in advance that a competitor is regular.

On $\mathcal V$ the recovery identity is $R\circ F=\mathrm{id}$.
Differentiation gives $DR\,J=I_p$, so $J$ is injective. The constant-rank
theorem identifies the local image as an embedded manifold, and the
finite free action gives its quotient charts. Internal root simplicity
makes the continuously labelled component roots analytic. Equality of
a root from one component with a root from another neither identifies
the complete tuples nor introduces a permutation stabilizer. It is
therefore immaterial to this recovery construction.
\end{proof}

Let $g(\theta)$ list the component roots, increasingly within each
component, and put $G=Dg(\theta_0)$ and $q_0=kd$. This labelling is a
local device, not an additional target supplied to an estimator.
Differentiating $f_b(\lambda_{ba})=0$ gives
\begin{equation}\label{eq:root-coefficient-derivative}
 \frac{\partial\lambda_{ba}}{\partial f_{bj}}
       =-\frac{\lambda_{ba}^{j}}{f_b'(\lambda_{ba})},\qquad j<d.
\end{equation}
For each component this is a nonsingular diagonal matrix times a
Vandermonde matrix. In particular $G$ has full row rank $q_0$.

\begin{theorem}[Two-sided local conditioning]\label{thm:exact-condition}
Let $H$ be a positive definite form on the product probability-simplex
tangent space
\[
 \mathcal T=\{Z:\ \one^{\mathsf T}\operatorname{vec}Z_j=0
                                      \text{ for every }j\}.
\]
It may be extended arbitrarily as a positive definite form to the
ambient space. Set
\begin{equation}\label{eq:intrinsic-covariance}
 V_H(P)=G(J^{\mathsf T}HJ)^{-1}G^{\mathsf T}.
\end{equation}
This positive definite $q_0\times q_0$ matrix is determined by the
observation and the specified metric, up to simultaneous permutation
of its rows and columns. The exact local inverse condition is
\begin{align}
 \mathfrak c_H(P)
 &:=\lim_{\epsilon\downarrow0}\sup_{\substack{P'\in F(\cK_d)\\
                     0<\|P'-P\|_H\le\epsilon}}
 \frac{d_\infty(\Lambda(P'),\Lambda(P))}{\|P'-P\|_H}
   =\max_{a\le q_0}\sqrt{[V_H(P)]_{aa}}.       \label{eq:exact-condition}
\end{align}
The maximum first-order ratio is attained along a smooth admissible
path. For $H=I$, $\mathfrak c_2(P)\le C_\nu/\eta_d(P)$ on the
within-component $\nu$-separated class. The same comparison holds for
a fixed uniformly equivalent metric.
\end{theorem}
\begin{proof}
Theorem~\ref{thm:quotient-chart} accounts for every nearby model datum.
Writing $v=\theta'-\theta_0$, Taylor expansion gives
$F(\theta')-P=Jv+O(\|v\|^2)$ and
$g(\theta')-g(\theta_0)=Gv+O(\|v\|^2)$.
Distinct values in the base aggregate multiset have a positive gap.
For sufficiently small $v$, matching cannot improve upon matching
within the corresponding base-value groups: a crossing between groups
has a fixed positive cost, whereas the within-group matching cost
tends to zero. Within an equal-value group every base root is the same
number. Hence its bottleneck distance from the perturbed roots is
exactly their largest absolute displacement. Thus, even at a collision,
\[
 d_\infty(\Lambda(\theta'),\Lambda(\theta_0))
       =\|Gv\|_\infty+O(\|v\|^2).
\]
This is a radial metric identity at the base datum; it does not assert
that sorting is differentiable. Injectivity of $J$ makes the desired
limit $\sup_{v\ne0}\|Gv\|_\infty/\|Jv\|_H$.
For each row of $G$, duality for the Gram matrix $J^{\mathsf T}HJ$
gives \eqref{eq:exact-condition}. A maximizing row and its dual vector
supply an admissible two-sided line in the open parameter chart.
The global comparison follows by differentiating
Theorem~\ref{thm:joint}. Under a coordinate change with invertible
Jacobian $D$, the matrices become $JD$ and $GD$; substitution leaves
\eqref{eq:intrinsic-covariance} unchanged. Component relabelling only
permutes its indices. Full row rank of $G$ proves positive definiteness.
Since columns of $J$ belong to $\mathcal T$, the ambient extension of
$H$ has no effect.
\end{proof}

For balanced categorical sampling the per-observation Fisher form on
$\mathcal T$ is
\[
 H_P=\ell^{-1}\diag\big(1/(P_j)_{ab}\big)_{j,a,b}.
\]
Section~\ref{sec:local-minimax} will identify the exact local statistical
meaning of the whole matrix $V_{H_P}$, including collision points.

An alternative derivative formula avoids differentiating an eigenvector
normalization. For an admissible observation tangent $Z$, let
\begin{align}
 r_Z&=\cN_P^\dagger\big[-(\Omega\otimes I)
                    (q(T_j)\operatorname{vec}Z_j)_j\big],
                                      \label{eq:normal-tangent}\\
 \dot M_Z&=Q_1^{\mathsf T}I_d(qZ+r_ZP)Q_2.\nonumber
\end{align}
The first expression denotes the polynomial with that coefficient
vector; the frames in the second are held fixed at the base point.
For an aggregate simple root $x$,
\begin{equation}\label{eq:root-tangent}
 \dot x=-\operatorname{tr}\big((\Res_{z=x}M(z)^{-1})\dot M_Z(x)\big).
\end{equation}
Differentiate $\cN_Pa=b_P$ for the first formula and the simple zero
of $\det M$ for the second. At a cross-component collision, one instead
differentiates the distinct joint tuples and their internally simple
roots using \eqref{eq:root-coefficient-derivative}. The trace of the
combined residue gives a sum of velocities, not the individual
velocities. The exact local condition is not being identified with a
distance to a discriminant or with the global certificate $\eta_d$.
